\section{Q1 --- with a perfect model, where does each scheme's error go?}\label{sec:q1}

\subsection{Overview}
Table~\ref{tab:overview} gives the RMSE of every scheme in both cases and on both
observing systems. With a perfect model, the learned schemes lead by 30 to
45\,\%:
\begin{itemize}[leftmargin=*,nosep]
  \item DirectUNet 0.340 and PredictStateCFM 0.341 on the regular grid;
  \item the best DA scheme, the ETKS, 0.497;
  \item Strong-4DVar, which assimilates with the true equations, 0.703.
\end{itemize}
The rest of this section asks where the DA deficit comes from.

\begin{table}[t]
\centering\small
\caption{RMSE on the 24 observed channels, 200 test windows; learned rows: 1200 epochs, 3 seeds pooled per window. The italic row is a sensitivity row. Sources: \texttt{reports/l96/outputs/p1\_l96\_benchmark.md}, \texttt{docs/results/l96\_weak4dvar.md}.}
\label{tab:overview}
\begin{tabular}{lcccc}
\toprule
 & \multicolumn{2}{c}{\Sz{} (perfect model)} & \multicolumn{2}{c}{\So{} (model error)} \\
\cmidrule(lr){2-3}\cmidrule(lr){4-5}
scheme & regular & random & regular & random \\
\midrule
ETKF & 0.610 & 0.679 & 1.409 & 1.407 \\
EnKF & 0.641 & 0.706 & 1.416 & 1.484 \\
ETKS & 0.497 & 0.572 & 1.338 & 1.347 \\
Strong-4DVar & 0.703 & 0.742 & 1.436 & 1.444 \\
Weak-4DVar & 0.695 & \needsrun{} & 1.064 & \needsrun{} \\
\emph{ETKS, 100 members} & \emph{0.393} & \emph{0.463} & \emph{1.338} & \emph{1.335} \\
\midrule
DirectUNet & 0.340 & 0.450 & 0.339 & 0.444 \\
PredictStateCFM & 0.341 & 0.443 & 0.338 & 0.447 \\
VanillaCFM & 0.351 & 0.462 & 0.349 & 0.468 \\
SDA1 & 0.464 & 0.572 & 0.462 & 0.585 \\
SDA2 & 0.453 & 0.550 & 0.465 & 0.575 \\
SDA3-fix & 0.455 & 0.562 & 0.455 & 0.573 \\
hybrid DirectUNet $\to$ SDA3-fix & \textbf{0.293} & \textbf{0.367} & \textbf{0.289} & \textbf{0.367} \\
\bottomrule
\end{tabular}
\end{table}

\subsection{The budget at \texorpdfstring{$\tau = 0$}{tau = 0}}\label{sec:q1-budget}

\paragraph{Irreducible.}
The best scheme, the hybrid, reaches 0.293 on the regular grid: the irreducible
error is at most this. We have no lower bound, so every share below is a share of
the excess over an upper bound, not over the true $\mathrm{mmse}$.

\paragraph{Restriction: filtering.}
The ETKS removes 16 to 19\,\% of the ETKF's RMSE (0.610 $\to$ 0.497 regular,
0.679 $\to$ 0.572 random). To check that this gain is a restriction, we run the
filter and the smoother on the same ensemble realisation and evaluate both sides
of the identity
$\mathrm{MSE}_F - \mathrm{MSE}_S = \mathrm{var}_F - \mathrm{var}_S = \E\|m_F - m_S\|^2$.
It holds approximately with 100 members --- an MSE gap of 0.162 against a
variance gap of 0.176 and a mean shift of 0.233 --- and less well with 30
(0.159, 0.227, 0.296). The deviation shrinks as the ensemble grows, as it should
if it is an approximation error of the ensemble.

\paragraph{Approximation: finite ensemble.}
With 100 members and a re-selected inflation, the ETKS goes from 0.497 to 0.393
(regular) and from 0.572 to 0.463 (random): 38\,\% of its mean-squared error was
finite-ensemble approximation. Against the best learned scheme, the DA deficit
falls from $1.46\times$ to $1.16\times$.

\paragraph{Approximation: optimisation and training.}
\begin{itemize}[leftmargin=*,nosep]
  \item \textbf{4D-Var.} At the benchmark's 10 L-BFGS iterations per sub-window,
    the hard-constraint 4D-Var is far from its own optimum: 0.721 against 0.545
    on the validation windows with 40 iterations.
  \item \textbf{DirectUNet.} Training it for 1200 instead of 400 epochs takes it
    from 0.380 to 0.340, and three times more training windows add only 0.006
    (0.334). The learned approximation term is close to saturation at our
    budget.
\end{itemize}

\paragraph{Restriction: context.}
The unconditional SDA prior (SDA1, 0.464) ignores $\thv$. Conditioning it on
$\thv$ gains 2\,\% (SDA2 0.453, SDA3-fix 0.455). \needsrun{The forcing $\zv$ as a
conditioning input of SDA, and $(\thv, \zv)$ as inputs of DirectUNet, complete
this row.}

\paragraph{Structural misspecification.}
What remains of the DA excess after these contrasts is structural: the mode
instead of the mean, a hard constraint and Gaussian updates. It is largest for
Strong-4DVar.
\todo{Table~5: the budget in pooled MSE per scheme class, both observing
systems, with the contrast named in each cell.}

\paragraph{Reading.}
With a perfect model, the deficit of classical DA is mostly \emph{honest}. It is
made of a restriction (filtering) and an approximation (30 members, a short
optimisation), and it shrinks with smoothing, with members and with iterations.
None of it requires a better model.

\subsection{The budget along \texorpdfstring{$\tau$}{tau}}
Table~\ref{tab:fms-s0} gives $\FMS_\tau$ for the main schemes at \Sz{} on the regular grid.
\begin{itemize}[leftmargin=*]
  \item \textbf{The smoother buys its mean with an over-confident ensemble.} It
    beats the filter at $\tau = 0$ (0.129 vs 0.187) and still at $\tau = 0.5$,
    but loses at $\tau = 0.75$ (0.059 vs 0.050). Its calibration loss there is
    34\,\% against 14\,\%.
  \item \textbf{The flows are nearly calibrated} (calibration loss 0--3\,\%,
    pooled spread/skill 0.84--0.94).
  \item \textbf{The SDA samplers are under-dispersed} (7--18\,\%), as expected
    from their tempered likelihood.
\end{itemize}
\needsrun{The filter/smoother identity along $\tau$ (score of the ETKS's own
filter pass) and the 100-member rows of this table.}

\begin{table}[t]
\centering\small
\caption{$\FMS_\tau$ (standardised units, Gaussian form) and calibration loss at $\tau = 0.75$, \Sz, regular grid. Source: \texttt{docs/results/l96\_p1\_fm\_score.md}.}
\label{tab:fms-s0}
\begin{tabular}{lcccc}
\toprule
scheme & $\tau = 0$ & $\tau = 0.5$ & $\tau = 0.75$ & calibration loss \\
\midrule
ETKF & 0.187 & 0.125 & 0.050 & 14\,\% \\
ETKS & 0.129 & 0.107 & 0.059 & 34\,\% \\
PredictStateCFM & 0.054 & 0.040 & \textbf{0.0195} & 2\,\% \\
VanillaCFM & 0.057 & 0.042 & 0.0196 & 0\,\% \\
SDA3-fix & 0.088 & 0.067 & 0.034 & 14\,\% \\
hybrid & \textbf{0.039} & \textbf{0.034} & 0.0200 & 16\,\% \\
\bottomrule
\end{tabular}
\end{table}

\figplaceholder{Fig.~3}{$\FMS_\tau$ curves per scheme, \Sz, regular and random observing systems, with window-bootstrap intervals.}

\subsection{Observing system and density}\label{sec:q1-density}
Figure~4 shows how the schemes depend on the number of observations.
\begin{itemize}[leftmargin=*]
  \item \textbf{Crossover.} DA is best when a window has at most about ten
    observation times; the learned schemes win from about twenty.
  \item \textbf{Transfer between observing systems.} The ratio of random to
    regular RMSE is 1.06 for Strong-4DVar, 1.15 for the ETKS, 1.21--1.24 for
    SDA and 1.30--1.32 for the CFMs and DirectUNet. Keeping $\Hop$ and $\mathbf{R}$
    explicit buys robustness to the observing system.
  \item \textbf{Beyond the training range} (from 300 to 1000 observations),
    DirectUNet degrades from 0.17 to 0.34, while SDA, whose likelihood is
    explicit, keeps improving from 0.30 to 0.25.
\end{itemize}
In the language of the chain, an amortised map is misspecified once the
observing system leaves its training distribution; a factorised scheme is not.

\figplaceholder{Fig.~4}{RMSE vs number of observation times and of observed fast channels, \Sz, all schemes.}
